\newpage
\section{Testowanie}	%5

\hspace{0.60cm}W trakcie wytwarzania oprogramowania prowadzono systematyczne testy funkcjonalne, sieciowe, interfejsu użytkownika (UI) oraz sprzętowe z wykorzystaniem emulatora systemu Android (Android Studio Emulator) oraz fizycznych urządzeń mobilnych. Celem weryfikacji było sprawdzenie stabilności odtwarzania strumieniowego, reakcji aplikacji na brak połączenia z internetem oraz poprawności działania czujników ruchu.

\subsection{Zestawienie scenariuszy i przypadków testowych}

\hspace{0.60cm}Poniższa tabela przedstawia zestawienie kluczowych scenariuszy testowych przeprowadzonych w ramach weryfikacji poprawności działania aplikacji „Guess the Beat”:

\begin{tabela}
	{Zestawienie przypadków testowych aplikacji mobilnej Guess the Beat}
	{Wyniki przeprowadzonych testów funkcjonalnych, sieciowych i sprzętowych}
	{
		\begin{tabular}{|c|p{3.5cm}|p{5.0cm}|p{3.0cm}|c|} \hline
			\textbf{ID} & \textbf{Obszar testowy} & \textbf{Opis scenariusza i dane wejściowe} & \textbf{Oczekiwany rezultat} & \textbf{Status} \\ \hline
			TC-01 & Wyszukiwanie playlist & Wpisanie frazy „Rock” w polu wyszukiwania na ekranie menu i kliknięcie przycisku szukania[cite: 12]. & Pobranie listy pasujących playlist z Deezer API i wyświetlenie ich w formie kart[cite: 8, 12]. & Pozytywny \\ \hline
			TC-02 & Obsługa braku internetu & Wyłączenie transmisji danych podczas ładowania playlisty lub utworu. & Przechwycenie wyjątku, powrót do menu głównego i wyświetlenie komunikatu o braku sieci[cite: 8, 10]. & Pozytywny \\ \hline
			TC-03 & Pomijanie utworów bez audio & Wylosowanie utworu z API, który posiada puste pole \texttt{preview}[cite: 8, 10]. & Automatyczne odrzucenie utworu z puli sesji i wylosowanie kolejnego bez błędu[cite: 10]. & Pozytywny \\ \hline
			TC-04 & Progresywne wydłużanie próbki & Kilkukrotne pominięcie etapu lub wpisanie błędnej odpowiedzi w rundzie[cite: 7, 10]. & Stopniowe wydłużanie czasu odsłuchu ($0,2\text{ s} \rightarrow 1\text{ s} \rightarrow 3\text{ s}$ itd.) i aktualizacja paska postępu[cite: 7, 9, 10]. & Pozytywny \\ \hline
			TC-05 & Lista podpowiedzi UI & Wpisanie min. 2 znaków w pole odpowiedzi przy wysuniętej klawiaturze ekranowej[cite: 9, 10]. & Wyświetlenie przewijanej listy podpowiedzi o ograniczonej wysokości (\texttt{max = 180.dp})[cite: 9]. & Pozytywny \\ \hline
			TC-06 & Cykl życia aktywności & Zminimalizowanie aplikacji lub odebranie połączenia podczas odtwarzania próbki[cite: 11]. & Natychmiastowe zatrzymanie dźwięku w metodzie \texttt{onPause()} i zwolnienie zasobów \texttt{MediaPlayer}[cite: 8, 11]. & Pozytywny \\ \hline
			TC-07 & Czujnik akcelerometru & Fizyczne lub symulowane odwrócenie telefonu ekranem do dołu w trakcie rundy[cite: 7]. & Jednorazowe wyzwolenie akcji pominięcia bieżącego etapu zgadywania[cite: 7]. & Wymaga poprawy \\ \hline
		\end{tabular}
	}
	\label{tab:testy}
\end{tabela}

\subsection{Testy interfejsu użytkownika i mechaniki etapów}

\hspace{0.60cm}Przeprowadzono szczegółowe testy manualne widoku \texttt{GameScreen.kt} pod kątem poprawności wyświetlania i ergonomii obsługi na różnych rozdzielczościach ekranu:
\begin{itemize}
    \item \textbf{Weryfikacja układu kontrolek:} Potwierdzono poprawne wyrównanie oraz jednakową wysokość pola tekstowego wyszukiwarki i przycisku pomijania etapu[cite: 7]. Sprawdzono również działanie dolnego paska nawigacyjnego (przejście do menu głównego oraz reakcję przycisku logowania)[cite: 7].
    \item \textbf{Testy autouzupełniania odpowiedzi:} Zweryfikowano poprawność działania funkcji \texttt{updateGuess}, która po wpisaniu minimum dwóch liter wysyła zapytanie do endpointu wyszukiwania utworów Deezer[cite: 8, 10]. Kliknięcie podpowiedzi poprawnie uzupełnia nazwę utworu i automatycznie zatwierdza odpowiedź[cite: 10].
    \item \textbf{Testy motywu kolorystycznego:} Sprawdzono działanie funkcji \texttt{isSystemInDarkTheme()} w \texttt{MainActivity.kt}, potwierdzając płynne przełączanie gradientu tła między trybem jasnym a ciemnym[cite: 11].
\end{itemize}

\subsection{Testy czujnika akcelerometru i analiza wykrytych problemów}

\hspace{0.60cm}Przetestowano działanie klasy \texttt{FlipDetector} zarówno przy użyciu wirtualnych czujników w emulatorze Android Studio (zakładka \textit{Extended Controls -> Virtual Sensors}), jak i poprzez fizyczne obracanie smartfona podczas rozgrywki[cite: 7]:
\begin{itemize}
    \item \textbf{Wykryty problem nadmiernej czułości (TC-07):} Podczas testów zaobserwowano, że czujnik akcelerometru reaguje zbyt gwałtownie na niewielkie przechylenia urządzenia[cite: 7]. Pojedynczy obrót telefonu potrafi kilkukrotnie pod rząd wyzwolić zdarzenie pominięcia etapu (\textit{skip}), co prowadzi do natychmiastowej utraty wszystkich prób w danej rundzie[cite: 7].
    \item \textbf{Przyczyna techniczna:} Zdarzenie \texttt{onSensorChanged} wywoływane jest z wysoką częstotliwością (kilkadziesiąt razy na sekundę), a brak programowej blokady czasowej (tzw. \textit{debouncingu}) oraz histerezy powrotu urządzenia do pozycji wyjściowej powoduje wielokrotne spełnienie warunku odwrócenia podczas jednego ruchu ręką[cite: 7].
    \item \textbf{Wnioski i planowane poprawki:} W ramach kolejnej iteracji prac zaplanowano kalibrację progów detekcji osi $Z$, zmniejszenie czułości czujnika oraz dodanie blokady czasowej uniemożliwiającej ponowne wywołanie funkcji pomijania przed upływem określonego czasu lub przed ponownym odwróceniem telefonu ekranem do góry[cite: 7].
\end{itemize}